\selectlanguage{ngerman}%
\frenchspacing%
\begin{ledgroupsized}[r]{120mm}%
\footnotesize%
\pstart%
\noindent\textbf{Überlieferung:}%
\pend%
\end{ledgroupsized}%
\begin{ledgroupsized}[r]{114mm}%
\footnotesize%
\pstart%
\parindent -6mm%
\makebox[6mm][l]{\textit{L}}%
Aufzeichnung:
LH~XXXV~10,\;14 Bl.\;12.
Ein Blatt 4\textsuperscript{o},
Ränder zum Teil beschnitten;
Fragment eines \href{\Scan{LH_35_10_14_012wz}}{Wasserzeichens}:
Wiener Papier.
Eine Seite auf Bl.\;\href{\Scan{LH_35_10_14_012r}}{12\,r\textsuperscript{o}} und sieben Zeilen auf Bl.\;\href{\Scan{LH_35_10_14_012v}}{12\,v\textsuperscript{o}}, sonst unbeschrieben.
Im Kopf von Bl.\;\href{\Scan{LH_35_10_14_012r}}{12\,r\textsuperscript{o}} ist von unklarer Hand die römische Zahl \textit{XI} in roter Tinte vermerkt
(vgl. die Träger der benachbarten Signaturen LH~XXXV~10,\;14 Bl.\;2 bis~13,
u.\,a.
N.\;\ref{60117},
\ref{60115},
\ref{60112},
\ref{60116},
\ref{60113},
\ref{60114},
\ref{60110}, 
\ref{41824}
u.\;\ref{60111}
im vorliegenden Band).
\pend%
\end{ledgroupsized}%
\vspace{5mm}%
\begin{ledgroup}%
\footnotesize%
\pstart%
\noindent%
\textbf{Datierungsgründe:}
Die Datierung der Aufzeichnung N.\;\ref{60118} beruht auf Eigenschaften des Textträgers.
Siehe die ausführliche Begründung auf S.\;\pageref{60118_dat}.%
\pend%
\end{ledgroup}%
\selectlanguage{latin}%
\frenchspacing%
\vspace{8mm}%
\normalsize%
\pstart%
\noindent%
\edlabel{LH_35_10_14_012_012r_anfg-1}%
\bracks{%
\href{\Scan{LH_35_10_14_012r}}%
{12\,r\textsuperscript{o}}%
}%
\edtext{}{%
\xxref{LH_35_10_14_012_012r_anfg-1}{LH_35_10_14_012_012r_anfg-2}%
\lemma{\bracks{\href{\Scan{LH_35_10_14_012r}}{12\,r\textsuperscript{o}}}}%
\Bfootnote{%
\textit{(1)}~Quaeri \bracks{/} %
\textit{(2)}~Quandonam%
~\textit{L}%
}} %
Quandonam%
\edlabel{LH_35_10_14_012_012r_anfg-2}
cubus%
\indexS{cubus gravis homogeneus}
\edtext{%
homogeneus gravis
cujus centrum \textit{C},%
\indexS{centrum cubi}%
}{%
\lemma{homogeneus gravis}%
\Bfootnote{%
\textit{(1)}~cujus hedra \angles{--} %
\textit{(2)}~cujus centrum \textit{C} %
\textit{erg.\;L}%
}}
in plano inclinato%
\indexS{planum inclinatum}
aliqua
\edtext{%
sua hedra \textit{AB} positus,%
\indexS{cubus in plano inclinato}%
\indexS{hedra cubi}%
}{%
\lemma{sua}%
\Bfootnote{%
\textit{(1)}~superfi %
\textit{(2)}~hedra \bars{\textit{AB} \textit{erg.}} positus,%
~\textit{L}%
}}
\edtext{%
gliscat%
\indexS{cubus gliscans}%
}{%
\lemma{gliscat}%
\Cfootnote{%
Wohl in der Bedeutung von \textit{labatur}, "gleitet". So auch bei den folgenden Verwendungen.%
}}
aut
\edlabel{LH_35_10_14_012_012r_fwc_1a}%
evertatur,%
\indexS{cubus eversus}%
\indexS{eversio cubi in plano inclinato}
et an
aliquando stet sine descensu%
\bracks{?}%
\indexS{descensus cubi in plano inclinato}%
\edtext{}{%
\xxref{LH_35_10_14_012_012r_fwc_1a}{LH_35_10_14_012_012r_fwc_1b}%
\lemma{evertatur,}%
\Bfootnote{%
\textit{(1)}~quaestio est. %
\bracks{/} Sit planum inclinatum \textit{AB}, centrum cubi \textit{C}, Horizon \textit{HB}. %
\textit{(2)}~et an %
\bracks{...} sine descensu \bracks{/} %
\textit{(a)}~Respondeo %
\textit{(aa)}~globum gliscere sine eversione %
\textit{(bb)}~gliscere hedra %
\textit{(b)}~Respondeo \textls{primo} %
\textit{(aa)}~globum\indexS{globus in plano inclinato} %
\textit{(bb)}~cubum semper descendere. Nam %
\textit{(c)}~Respondebo per aliquot positiones: \bracks{/} %
\textit{(aa)}~\textls{Primo.} Globus quamdiu gliscit, descendit %
\textit{(bb)}~\textls{Primo.} %
\textit{(aaa)}~E\angles{--} %
\textit{(bbb)}~Eversio fit %
\bracks{...} fulcro nempe %
\bars{super \textit{erg.}} ima acie hedrae %
\bars{cujus unum %
\bars{medium \textit{erg., ändert Hrsg.}} punctum \textit{B} hic designo \textit{erg.}}.%
~\textit{L}%
}} %
Respondebo per aliquot positiones:%
\indexS{positio in responsione}
\pend%
\begin{edfigure}[2.25em][0.00em]
\hspace{-84mm}%
\includegraphics%
{gesamttex/edit_VIII-4/images/LH_35_10_14_012_d1_012r.pdf}%
\edcaption[0.75em]{%
\hspace{-84mm}%
\bracks{\textit{Fig.\:1}}%
}%
\vspace{-8.675em}%
\includegraphics%
{gesamttex/edit_VIII-4/images/LH_35_10_14_012_d2_012r.pdf}%
\edcaption[0.675em]{%
\bracks{\textit{Fig.\:2}}%
}%
\vspace{-8.75em}%
\hspace{84mm}%
\includegraphics%
{gesamttex/edit_VIII-4/images/LH_35_10_14_012_d3_012r.pdf}%
\edcaption[0.525em]{%
\hspace{84mm}%
\bracks{\textit{Fig.\:3}}%
\edtext{}{%
\lemma{\bracks{\textit{Fig.\:1}} bis \bracks{\textit{Fig.\:3}}}%
\Cfootnote[nonum]{%
Zusätzlich zu den Diagrammen überliefert \textit{L} (Bl.\;\href{\Scan{LH_35_10_14_012r}}{12\,r\textsuperscript{o}}) zwei gestrichene, hier nicht wiedergegebene Entwürfe.%
}}
}%
\end{edfigure}
\newpage%
\pstart%
\textls{Primo.}
Eversio%
\indexS{eversio cubi in plano inclinato}
fit super \textit{B} tanquam fulcro%
\indexS{fulcrum eversionis cubi}%
\bracks{,}
nempe super ima acie hedrae%
\indexS{hedra cubi}%
\indexS{acies hedrae}
cujus unum \bracks{extrem}um punctum \textit{B} hic designo.%
\edlabel{LH_35_10_14_012_012r_fwc_1b}
\pend%
\pstart%
\textls{Secundo.}
Itaque
si recta \textit{CB} %
sit horizonti \textit{HB}%
\indexS{horizon}
\edtext{%
perpendicularis%
\indexS{recta perpendicularis horizonti}
\edtext{%
(minima a centro%
\indexS{centrum cubi}
ad fulcrum)%
\indexS{fulcrum eversionis cubi}%
}{%
\lemma{(minima \bracks{...} fulcrum)}%
\Cfootnote{%
Das zur Geraden \textit{HB} senkrechte Segment \textit{CB} stellt vielmehr den größten Abstand des Zentrums des Würfels vom Horizont dar.%
}}
(ut%
}{%
\lemma{perpendicularis}%
\Bfootnote{%
\bars{%
\textit{(1)}~minima %
\textit{(2)}~(minima a centro ad fulcrum) \textit{erg.}} (ut%
~\textit{L}%
}}
in
\edtext{%
fig.\;2)%
\indexS{figura}%
}{%
\lemma{fig.\;2}%
\Cfootnote{%
Vgl. das Diagramm \bracks{\textit{Fig.\;2}}.%
}}
aequilibrium est super%
\indexS{aequilibrium cubi super fulcro} %
hoc
\edlabel{LH_35_10_14_012_012r_dbn-1}%
fulcro,%
\indexS{fulcrum eversionis cubi}%
\edtext{}{%
\xxref{LH_35_10_14_012_012r_dbn-1}{LH_35_10_14_012_012r_dbn-2}%
\lemma{fulcro}%
\Bfootnote{%
\textit{(1)}~. \bracks{/} \textls{Tertio} %
\textit{(2)}~, nam %
\textit{(a)}~\textit{CB} %
\textit{(b)}~planum verticale %
\bracks{...} partes similes %
\textit{(aa)}~aequales, %
\textit{(bb)}~et aequales, %
\bracks{...} sine eversione. \bracks{/} \textls{Tertio.}%
~\textit{L}%
}}
nam planum verticale
per \textit{CB} et aciem \textit{B}
dividit cubum%
\indexS{cubus in plano inclinato}
in duas partes
similes et aequales,
et quantum pars dextra
(in figura)%
\indexS{figura}
nititur ad eversionem,%
\indexS{nisus ad eversionem}
tantum sinistra contra;
manet igitur cubus adhuc sine eversione.%
\indexS{eversio cubi in plano inclinato}
\pend%
\pstart%
\textls{Tertio.}%
\edlabel{LH_35_10_14_012_012r_dbn-2}
Itaque hic status%
\indexS{status cubi in plano inclinato}
est terminus gliscendi,%
\indexS{terminus gliscendi}
hoc est elevatio plani inclinati semirecta%
\indexS{elevatio plani inclinati}
est maxima earum in quibus
\edtext{%
\bracks{cubus}%
}{%
\lemma{cubus}%
\Bfootnote{%
\textit{erg. Hrsg.}%
}}
gliscitur sine eversione,%
\indexS{cubus gliscans}
nam si angulus \textit{CBH} est rectus,
ideo cum angulus \textit{CBA} sit semirectus,
residuus \textit{ABH} erit etiam semirectus.
\pend%
\pstart%
\textls{Quarto.}
Si major
\edtext{%
sit elevatio quam semirecta,%
\indexS{elevatio plani inclinati}
cubus%
}{%
\lemma{sit}%
\Bfootnote{%
\textit{(1)}~inclinatio %
\textit{(2)}~elevatio \bars{ quam \textit{erg.}} semirecta, %
\textit{(a)}~globus %
\textit{(b)}~cubus%
~\textit{L}%
}}
evertitur%
\indexS{cubus eversus}%
\indexS{eversio cubi in plano inclinato}
ut in
\edtext{%
fig.\;3.%
\indexS{figura}%
}{%
\lemma{fig.\;3}%
\Cfootnote{%
Vgl. das Diagramm \bracks{\textit{Fig.\;3}}.%
}}
Exceditur enim aequilibrium%
\indexS{aequilibrium cubi super fulcro}
casus praecedentis.%
\indexS{casus aequilibrii}
\pend%
\pstart%
\textls{Quinto.}
Si minor
\edtext{sit,}{%
\lemma{sit}%
\Bfootnote{%
\textit{erg.\;L}%
}}
gliscit tantum,%
\indexS{cubus gliscans}
ut patet.
\pend%
\pstart%
\textls{Sexto.}
\edtext{%
Quamdiu cubus gliscere potest%
\indexS{cubus gliscans}
descendit,%
\indexS{cubus descendens}%
\indexS{descensus cubi in plano inclinato}%
}{%
\lemma{Quamdiu}%
\Bfootnote{%
\textit{(1)}~globus %
\textit{(2)}~cubus %
\textit{(a)}~gliscit descendit %
\textit{(b)}~gliscere potest descendit,%
~\textit{L}%
}}
potest enim centrum gravitatis descendere%
\indexS{centrum gravitatis cubi}%
\indexS{descensus centri gravitatis}
in linea recta
\edtext{%
plano inclinato%
\indexS{planum inclinatum}
parallela.%
}{%
\lemma{plano}%
\Bfootnote{%
\textit{(1)}~parallela. %
\textit{(2)}~inclinato parallela.%
~\textit{L}%
}}
\pend%
\pstart%
\textls{Septimo.}
Cum cubus evertitur%
\indexS{cubus eversus}%
\indexS{eversio cubi in plano inclinato}
descendit;%
\indexS{descensus cubi in plano inclinato}
nam si descendit cubus%
\indexS{cubus descendens}
per
\edtext{%
planum%
\indexS{planum inclinatum}
ut cum gliscit%
\indexS{cubus gliscans}
per art.\;6,%
\indexS{articulus}
id est in minore elevatione%
\indexS{elevatio plani inclinati}
per artic.\;3,
multo magis descendet in majore elevatione,%
}{%
\lemma{planum}%
\Bfootnote{%
\textit{(1)}~in minore elevatione, multo magis ascendet in majore elevati %
\textit{(2)}~cum gliscit %
\textit{(3)}~ut cum \bracks{...} majore elevatione,%
~\textit{L}%
}}
hoc est per artic.\;4,%
\indexS{articulus}
cum evertitur.%
\indexS{cubus eversus}%
\indexS{eversio cubi in plano inclinato}
\pend%
\pstart%
\textls{Octavo.}
Itaque
\edtext{%
cubus homogeneus%
\indexS{cubus gravis homogeneus}
in plano%
}{%
\lemma{cubus}%
\Bfootnote{%
\textit{(1)}~in pl %
\textit{(2)}~homo %
\textit{(3)}~homogeneus in plano%
~\textit{L}%
}}
inclinato%
\indexS{planum inclinatum}
positus%
\indexS{cubus in plano inclinato}
semper descendit.%
\indexS{cubus descendens}
Nam descendit%
\indexS{descensus cubi in plano inclinato}
cum 
\edtext{plani}{%
\lemma{plani}%
\Bfootnote{%
\textit{erg.\;L}%
}}
elevatio est%
\indexS{elevatio plani inclinati}
\edtext{%
infra semirectam,
quia tunc per 3 aut 5 gliscit,%
\indexS{cubus gliscans}
ergo per 6 descendit.%
}{%
\lemma{infra}%
\Bfootnote{%
\textit{(1)}~grad %
\textit{(2)}~semirectam, %
\textit{(a)}~per 6 et 3 et 5 %
\textit{(b)}~quia tunc %
\bracks{...} per 6 descendit.%
~\textit{L}%
}}
Cum vero plani elevatio%
\indexS{elevatio plani inclinati}
est major,
descendit multo
magis.%
\indexS{cubus descendens}%
\indexS{descensus cubi in plano inclinato}
\bracks{%
\href{\Scan{LH_35_10_14_012v}}%
{12\,v\textsuperscript{o}}%
}
\pend%
\pstart%
\textls{Nono.}
\ledpb%
In eversione compositus est motus,%
\indexS{motus compositus ex eversione et descensu}
ex eversione %
(quae eodem modo fieret et in plano horizontali \textit{HB})%
\indexS{planum horizontale}
et ex descensu,%
\indexS{descensus cubi in plano inclinato}
qui interim
cum fieri possit,
ab eversione%
\indexS{eversio cubi in plano inclinato}
non retardatur.
\pend%
\pstart%
\textls{Decimo.}
Operae pretium ergo foret definire lineam,
quam describit centrum cubi%
\indexS{centrum cubi eversi}
durante motu composito ex eversione et descensu.%
\indexS{motus compositus ex eversione et descensu}
\pend
\vspace{5.0em}%
\pstart%
\noindent%
\centering%
*     *     *
\pend%
\vspace{5.0em}%
\selectlanguage{german}%
\frenchspacing%
\footnotesize%
\pstart%
\noindent%
\textls{Ausführliche Datierungsgründe:}%
\pagelabel{60118_dat}
Der Träger der titellosen und undatierten Aufzeichnung N.\;\ref{60118} weist das Fragment eines \href{\Scan{LH_35_10_14_012wz}}{Wasserzeichens} auf, das die Wiener\indexO{Wien} Herkunft des Papiers bezeugt. %
Nach heutigem Wissensstand ist dieses Papier nur für den ersten und den zweiten Aufenthalt Leibnizens in der habsburgischen Hauptstadt\indexO{Wien} belegt, d.\,h. für den Zeitraum vom 8.\;Mai 1688 bis etwa zum 12.\;Februar 1689 sowie für den Zeitraum von Ende April bis Mitte Mai 1690 (vgl. \textit{Chronik}, S.\;89--95 u.\;103).\cite{01236} 
Dass in den Monaten dazwischen Leibniz auf seiner Italienreise\indexO{Italien} mitgenommenes Papier dieser Art gelegentlich verwendet haben kann, ist zwar unwahrscheinlich, aber nicht unmöglich.
Hieraus ergeben sich insgesamt der früheste und der späteste Entstehungszeitpunkt der Aufzeichnung N.\;\ref{60118}.
\pend%
\pstart%
Während ein früheres Entstehen derselben %
so gut wie ausgeschlossen ist, ist eine spätere Datierung nicht undenkbar, war doch Leibniz zeit seines Lebens noch dreimal in Wien,\indexO{Wien} nämlich von Ende Oktober bis Mitte Dezember 1700, im Dezember 1708 und von Mitte Dezember 1712 bis Anfang September 1714 (vgl. \textit{Chronik}, S.\;167--169, 211\,f. u.\;233--247).\cite{01236}
Die Aufzeichnung N.\;\ref{60118} könnte auch bei einem dieser späteren Aufenthalte verfasst worden sein.
Für eine Verwendung des besagten Papiers nach 1690 ist im Leibniz-Nachlass jedoch kein Beleg bekannt.
\pend%
\selectlanguage{latin}%
\frenchspacing%
\normalsize%
